\section{Conclusion}
\label{sec:conclusion}

We set out to provide a principled, measurable decision framework for shopping
agents under hard resource constraints: determining when adaptive sequential search
delivers value over fixed stopping rules.

The answer has three parts. Offers in the live corpus are genuinely ephemeral, so
stopping is a real decision rather than a formality. An adaptive reservation-price
rule does beat the strongest tuned fixed rule on held-out live panels, with the whole
paired interval below zero and no budget violations. And the margin is modest in our
current snapshot, for a reason that is practically informative: when merchants charge
nearly identical prices, there is little price variation for an optimizer to exploit,
whereas calibrated simulations show substantial gains when dispersion and search
horizons widen.

That is why we report an explicit decision map rather than a single aggregate
headline. Mapping where adaptive search yields significant advantage versus where
simpler heuristics suffice provides actionable guidance for deploying autonomous
agents in live commerce environments.

The narrowness of the current two-snapshot benchmark is itself an empirical baseline.
Cross-merchant product overlap in this early ecosystem is thin, and half of the
observed episodes exhibit minimal price spread. Our framework establishes a rigorous
foundation to evaluate when adaptive search pays as multi-merchant catalog networks
expand.

We release the optimizer, the permit ledger, the replay harness, the collection
scripts, and the dated snapshots, so the measurement can be repeated on wider corpora
and longer horizons.
